\chapter{Implementation notes}\label{app-impl}

This appendix collects the differences between the implementations that
the book mentions: ga68, the Algol 68 front end of GCC; a68j, which
compiles the same language to the Java virtual machine; and Algol 68
Genie (a68g), an interpreter that many Algol 68 programs on the web,
among them most Rosetta Code programs, are written for. The information
describes the versions used for this book, which the preface lists.

\section*{Notation}

\begin{center}
\small
\begin{tabular}{p{0.2\textwidth}p{0.35\textwidth}p{0.35\textwidth}}
\toprule
 & ga68 and a68j & Algol 68 Genie \\
\midrule
reserved words & lower case: \code|begin int| & upper case: \texttt{BEGIN INT} \\
declared type names & capitalized: \code|Node| & upper case: \texttt{NODE} \\
word operators & upper case: \code|ABS| & upper case: \texttt{ABS} \\
names & lower case, underscores: \code|max_int| & lower case, blanks allowed: \texttt{max int} \\
comments & \code|{ ... }| (nested) & \texttt{\# ... \#}, \texttt{CO ... CO}, \texttt{COMMENT ... COMMENT} \\
program & a block or other enclosed construct & the same, or a sequence without brackets \\
\bottomrule
\end{tabular}
\end{center}

Converting a program from one notation to the other is mechanical: the
reserved words and the type names change case, blanks in names become
underscores, and comments get the other brackets. A name that is a
reserved word in ga68 but not in Algol 68 Genie, such as \code|in|,
\code|out|, \code|empty|, \code|word| or \code|op|, must be renamed.

\section*{Sizes}

\begin{center}
\begin{tabular}{lcc}
\toprule
 & ga68 and a68j & Algol 68 Genie \\
\midrule
\code|int| & 32 bits & 64 bits \\
\code|long int| & 64 bits & 128 bits \\
\code|long long int| & 64 bits & 72 decimal digits \\
\code|real| & 32 bits & 64 bits \\
\code|long real| & 64 bits & 128 bits \\
\code|long long real| & 80 bits (ga68), 64 bits (a68j) & 63 decimal digits \\
\code|bits| & 32 & 64 \\
\bottomrule
\end{tabular}
\end{center}

A program written for Algol 68 Genie that relies on 64-bit \code|int|
values must use \code|long int| with ga68 and a68j. Since \code|real| has
only about seven decimal digits in ga68 and a68j, numerical programs
should use \code|long real| there.

\section*{Transput}

ga68 and a68j have no transput of the Report in their standard library.
The module \code|Genie_Transput|, which comes with a68j and can be
compiled with ga68, provides unformatted transput: \code|print|,
\code|read|, \code|put|, \code|get|, files on disk, files associated with
strings, the layout procedures and the events for the end of the file and
for value errors. It does not provide formats, binary transput, or the
printing of records, \code|compl|, \code|short| and \code|long long|
values. Algol 68 Genie provides the transput of the Report almost
completely, and adds \code|eof| and \code|eoln|, which
\code|Genie_Transput| also has; Section~\ref{sec-transput-genie} lists
what it lacks.

\section*{Other features}

\begin{longtable}{p{0.3\textwidth}p{0.3\textwidth}p{0.3\textwidth}}
\toprule
Feature & ga68 and a68j & Algol 68 Genie \\
\midrule
\endhead
\code|par| & ga68: one unit after the other; a68j: threads with the option \code|-fpar-threads| & threads \\
\code|sema| & not in ga68; a68j with \code|-fpar-threads| & yes \\
\code|compl| operators & a68j only & yes \\
\code|bytes| & no & yes \\
formats & no & yes \\
lifetime rule (references, procedures, \code|loc| cells) & not checked & checked at run time \\
bounds of an array field when a record is assigned & not checked & checked \\
initial value of a variable & zero, \code|false|, blank, empty & undefined (an error when used) \\
\code|andth|, \code|orel| & yes (GNU extension) & yes (also \texttt{ANDF}, \texttt{ORF}) \\
\code|XOR| for \code|bits| & yes & yes \\
\code|downto|, \code|until| in loops & no & yes \\
modules & yes (GNU extension, Appendix~\ref{app-modules}) & no \\
calling C (\code|nest C|) & ga68 only (Appendix~\ref{app-nest}) & no \\
partial parametrization & no & yes (Appendix~\ref{app-partial}) \\
\bottomrule
\end{longtable}

\section*{Printing numbers}

\code|print| writes an \code|int| in a field of \code|int_width + 1|
characters with its sign, so the output of ga68 and a68j is narrower than
that of Algol 68 Genie. A \code|real| is written in floating-point
notation with \code|real_width| digits. Programs whose output should not
depend on the implementation use \code|whole| and \code|fixed|.

When ga68 and a68j convert a \code|real| to a string, the last of its
seven digits is sometimes wrong (9.81 is printed as \texttt{9.810001});
the examples of this book avoid such values.

\section*{Running Rosetta Code programs}

Most Algol 68 programs on Rosetta Code are written for Algol 68 Genie in
upper-case notation. To run one with ga68 or a68j, convert the notation as
above, add \code|access Genie_Transput| and an enclosing block if the
program has none, and check whether it relies on 64-bit integers, on
formats or on the other features in the table above.
